\scanpage{8}

\deleted{En vertu du th. 1 \deleted{\illegible{1}}, il opère sans pt fixes, donc le quotient}

\uline{Si} $F'$ est représentable par \inserted{un schéma \unsure{qu.|quasi} projectif sur $S$} $M'$, alors $A$ opère sur $M'$, en vertu du th. 1 il opère sans pts fixes. \margin{à dégager} D'après un critère général \inserted{$A$ étant plat sur $S$}, il s'ensuit que $F$ est représentable \unsure{ssi|car} $M'/A$ existe ``raisonnablement'', et alors $M=M'/A$ représente $F$. Or on sait passer au quotient puisque $A$ est \uline{propre} et plat sur $S$,

Or $F'(T)$ s'explicite comme sous-ensemble de $\underline{\mathrm{Hilb}}_{A/S}(T)$, puisque les $C$ qu'on veut classifier \unsure{équivalent|équivalentes} aux $C$ \uline{plongés dans} $\deleted{\unsure{C_T}}A_T$, tels que le morphisme naturel $J^0_{C/S}\longrightarrow A_T$ déduit du \unsure{plongemt|plongement} soit un \uline{isom. de schémas Abéliens polarisés}, \deleted{\unsure{i.e. le degré est où la polarisation}} \deleted{Or, \unsure{le|la} exigeons que} Ceci implique d'ailleurs que, dans le morphisme $C\longrightarrow A_T$, \deleted{le sous-schéma fermé $C$ de $A_T$ soit} la polarisation induite sur $C$ soit de degré $g$, donc donné. D'après la théorie du schéma de Hilbert, si on fait abstraction du fait que $J^0_{C/S}\to A_T$ doive être un

% archive sheet number: 7
